\medskip
\hypertarget{native-108-article079}{}\pdfbookmark[2]{79.}{article-079}
\begin{center}79.\end{center}
Ad completam determinationem sectionis conicae in plano suo \emph{tria} requiruntur, situs perihelii, excentricitas et semiparameter. Quae si e quantitatibus datis ab ipsis pendentibus eruenda sunt, tot data adsint oportet, ut tres aequationes ab invicem independentes formare liceat. Quilibet radius vector magnitudine et positione datus unam aequationem suppeditat: quamobrem ad determinationem orbitae tres radii vectores magnitudine et positione dati requiruntur; si vero duo tantum habentur, vel unum elementum ipsum iam datum esse debet, vel saltem alia quaedam quantitas, cui aequationem tertiam superstruere licet. Hinc oritur varietas problematum, quae iam deinceps pertractabimus.
\par
Sint $r$, $r'$ duo radii vectores, qui cum recta in plano orbitae e Sole ad lubitum ducta faciant secundum directionem motus angulos $N$, $N'$; sit porro \SourceLeaf{109}$\Pi$ angulus quem cum eadem recta facit radius vector in perihelio, ita ut radiis vectoribus $r$, $r'$ respondeant anomaliae verae $N-\Pi$, $N'-\Pi$; denique sit $e$ excentricitas, $p$ semiparameter. Tunc habentur aequationes
\[
\begin{aligned}
\frac{p}{r}&=1+e\cos(N-\Pi)\\
\frac{p}{r'}&=1+e\cos(N'-\Pi),
\end{aligned}
\]
e quibus, si insuper una quantitatum $p$, $e$, $\Pi$ data est, duas reliquas determinare licebit.
\par
Supponamus primo, datum esse semiparametrum $p$, patetque determinationem quantitatum $e$ et $\Pi$ ex aequationibus
\[
\begin{aligned}
e\cos(N-\Pi)&=\frac{p}{r}-1\\
e\cos(N'-\Pi)&=\frac{p}{r'}-1
\end{aligned}
\]
fieri posse ad normam lemmatis III in art. praec. Habemus itaque
\[
\begin{aligned}
\operatorname{tang}(N-\Pi)&=\operatorname{cotang}(N'-N)-\frac{r(p-r')}{r'(p-r)\sin(N'-N)}\\[5pt]
\operatorname{tang}(\tfrac12N+\tfrac12N'-\Pi)&=\frac{(r'-r)\operatorname{cotang}\tfrac12(N'-N)}{r'+r-\dfrac{2rr'}{p}}\mathpunct{\cdot}
\end{aligned}
\]
